% RecSys Odyssey repository-backed lecture note.
% Add figures with: \coursefigure{figures/name.svg}{Caption}
\section{Quality has more than one axis}

\begin{abstract}
A relevant slate can still be repetitive, narrow or unsurprising.
\end{abstract}

\subsection{Concept}
Intra-list diversity measures difference among items shown together. Coverage measures how much catalogue or category breadth a policy reaches across requests. Novelty rewards items unfamiliar to the user, while serendipity asks for useful surprise rather than obscurity for its own sake. These metrics should be read beside relevance: maximizing any one alone creates either a filter bubble or low-value filler.

\begin{equation*}
slate quality = relevance + diversity + novelty + constraints
\end{equation*}

\subsection{Key terms}
\begin{itemize}
\item \textbf{Intra-list diversity.} How dissimilar items within one slate are from each other.
\item \textbf{Novelty.} How unfamiliar a recommended item is to the user.
\item \textbf{Serendipity.} A relevant recommendation that is also meaningfully unexpected.
\end{itemize}
